\section{Conclusiones}

\subsection{Logros del trabajo}

Se cumplió el objetivo general del caso de estudio: se implementaron dos versiones paralelas del programa de multiplicación de matrices, una con \textbf{POSIX Threads} y otra con \textbf{procesos + memoria compartida POSIX}, y se compararon con la versión secuencial usando la métrica estándar de speedup.

El diseño experimental fue riguroso: 10 corridas por configuración, 5 tamaños de matriz, hasta 5 cantidades de hilos/procesos, en total 450 mediciones. El script se diseñó de forma que se puede pausar y reanudar sin perder datos, lo cual fue clave para un experimento de varias horas.

\subsection{Hallazgos principales}

\begin{enumerate}
    \item \textbf{No hay una técnica universalmente mejor}: la elección entre pthreads y fork depende del tamaño del problema. En este equipo, para matrices chicas ($N < 2000$) fork fue mejor; para matrices grandes ($N \geq 4000$) pthreads fue mejor.
    \item \textbf{El speedup real está lejos del ideal}. El mejor caso fue 6.90$\times$ con 8 hilos en $N = 8000$, mientras que el ideal sería 8$\times$ con 8 hilos o 12$\times$ con 12 procesos (recordando que la máquina tiene 12 núcleos lógicos). La diferencia se debe a la parte secuencial del programa, la contención por memoria y el \textit{overhead} del sistema.
    \item \textbf{Repetir el experimento 10 veces y usar la mediana fue clave}: se detectaron outliers puntuales debidos a ruido del sistema, y la mediana los ignoró de forma natural.
    \item \textbf{El orden de los bucles en el script}: alternar tamaños entre repeticiones evitó que los datos se quedaran ``calientes'' en la caché del procesador, dando mediciones más realistas.
\end{enumerate}

\subsection{Limitaciones}

\begin{itemize}
    \item El speedup de fork se calculó contra el $T = 1$ de \textbf{pthreads} (no se midió fork $T = 1$ por separado). En teoría fork $T = 1$ debería dar un tiempo similar al secuencial de pthreads, pero no se verificó experimentalmente.
    \item La versión fork con $N = 8000$ usa mucha memoria: tres matrices de 256\,MB cada una, más la imagen del proceso. En máquinas con menos RAM podría fallar.
    \item El experimento se hizo en un solo equipo. En otra máquina con distinto hardware (más o menos núcleos, diferente jerarquía de caché) los resultados podrían cambiar.
\end{itemize}

\subsection{Trabajo futuro}

Como continuación natural de este trabajo se podrían explorar:

\begin{itemize}
    \item Implementar la versión con procesos midiendo también \texttt{T=1}, para tener un baseline propio de fork.
    \item Probar tamaños de matriz aún más grandes ($N = 16000$ o más) para ver si la tendencia de pthreads ganando se mantiene.
    \item Aplicar optimizaciones al algoritmo (caché blocking, orden de los bucles cambiado) y medir cómo afectan al speedup.
    \item Probar en una máquina con muchos más núcleos (32 o 64) para ver hasta dónde escala cada técnica.
\end{itemize}
